%----------------------------------------------------------------------------- %%
\section{Related Work}
\label{bab:2}
%----------------------------------------------------------------------------- %%

Kubernetes distributes Pods across worker Nodes based on declared resource requests. These decisions are made when a Pod is scheduled, but because running Pods are not continuously relocated, an initially acceptable layout can become fragmented over time after scale-outs, terminations, and replacement events. To address this, the Kubernetes Descheduler evicts Pods that contribute to an undesirable placement state, allowing the kube-scheduler to place them again~\cite{k8s_descheduler_docs}. Upstream utilization policies, such as \textit{HighNodeUtilization}, pack workloads by draining low-utilization Nodes toward more utilized Nodes. However, these default strategies evaluate node capacity in isolation and do not account for the network relationships between the Pods they are evicting. Moving a Pod to satisfy a compute-utilization threshold can inadvertently sever its proximity to dependent services, leading to latency degradation.

In parallel, communication-aware Kubernetes studies recognize that microservices exchange traffic across service dependencies, and that placement across distant zones or degraded links increases latency even when compute resources fit perfectly. These studies incorporate topology, traffic, or measured network conditions into placement decisions~\cite{marchese2022network,marchese_tomarchio_2022,chen2024trade,pusztai_pogonip_2021}. For example, \cite{marchese2022network} and \cite{chen2024trade} demonstrate how placing highly communicative Pods on the same Node or within the same availability zone drastically reduces tail latencies. However, these works primarily target the initial scheduler or rely on atomic group migration. Moving dependent services as an atomic group can increase the risk of resource deadlock for a consolidation-oriented policy, as the scheduler must find contiguous capacity for a larger combined group.

This research bridges compute consolidation with network safety by applying network awareness as a pre-eviction guardrail rather than an initial placement scheduler. A global pre-eviction filter compares the current communication cost with feasible alternatives, intercepting harmful evictions proposed by standard compute-balancing policies. This allows standard consolidation policies to operate while providing a safety filter that prevents them from inadvertently destroying communication locality.
